\documentclass[11pt,a4paper]{article}
\usepackage[margin=25mm]{geometry}
\usepackage[T1]{fontenc}
\usepackage{lmodern,amsmath,amssymb,amsthm,mathtools,booktabs,array}
\usepackage{xcolor}
\usepackage[colorlinks=true,linkcolor=blue!45!black,urlcolor=blue!45!black,citecolor=blue!45!black]{hyperref}
\usepackage{microtype,fancyhdr}
\hypersetup{pdftitle={Grouped Diagonals and Arithmetic Bounds in Creative Telescoping},pdfauthor={Keio}}
\newtheorem{theorem}{Theorem}[section]
\newtheorem{lemma}[theorem]{Lemma}
\newtheorem{proposition}[theorem]{Proposition}
\newtheorem{corollary}[theorem]{Corollary}
\theoremstyle{remark}\newtheorem{remark}[theorem]{Remark}
\setlength{\emergencystretch}{3em}
\setlength{\headheight}{14pt}
\pagestyle{fancy}\fancyhf{}
\fancyhead[L]{\small Keio}
\fancyhead[R]{\small Grouped diagonals and arithmetic bounds}
\fancyfoot[C]{\thepage}
\title{Grouped Diagonals and Arithmetic Bounds\\in Creative Telescoping\\[5pt]
\large Constructive proofs addressing Chen--Kauers Problems 12 and 4}
\author{Keio\thanks{Research development and manuscript preparation with ChatGPT.}}
\date{2 November 2026}
\begin{filecontents*}[overwrite]{verify_grouped_diagonals.py}
"""Exact finite checks of the coefficient-splitting construction.

These checks supplement, and do not replace, the all-index proof.
Requires SymPy. Run with Python 3.
"""
import itertools
import json
from functools import lru_cache
from pathlib import Path
import sympy as S


def rational_coefficients(expr, variables):
    numerator, denominator = S.cancel(expr).as_numer_denom()
    p = S.Poly(numerator, *variables).as_dict()
    q = S.Poly(denominator, *variables).as_dict()
    zero = (0,) * len(variables)
    assert q.get(zero, 0) != 0, 'denominator must be a unit at the origin'
    q0 = q.pop(zero)

    @lru_cache(None)
    def coefficient(index):
        if any(a < 0 for a in index):
            return S.Integer(0)
        total = p.get(index, 0)
        for exponent, value in q.items():
            if all(a >= b for a, b in zip(index, exponent)):
                total -= value * coefficient(tuple(a-b for a,b in zip(index, exponent)))
        return total / q0
    return coefficient


def split(expr, old, left, right):
    return S.cancel((left * expr.subs(old, left)
                     - right * expr.subs(old, right)) / (left-right))


def construct(expr, old_variables, matrix):
    s = len(matrix[0])
    groups = [[] for _ in range(s)]
    out = expr
    for i, (old, row) in enumerate(zip(old_variables, matrix)):
        allocation = [j for j, count in enumerate(row) for _ in range(count)]
        if not allocation:
            out = out.subs(old, 0)
            continue
        variables = S.symbols('v%d_0:%d' % (i, len(allocation)))
        if len(variables) == 1:
            out = out.subs(old, variables[0])
        else:
            current = old
            for k in range(len(variables)-1):
                next_var = variables[-1] if k == len(variables)-2 else S.Dummy()
                out = split(out, current, variables[k], next_var)
                current = next_var
        for variable, j in zip(variables, allocation):
            groups[j].append(variable)
    for j, group in enumerate(groups):
        if not group:
            variable = S.Symbol('dummy%d' % j)
            group.append(variable)
            out /= 1-variable
    return S.cancel(out), groups


def main():
    z, w = S.symbols('z w')
    cases = [
        ('overlapping', 1/(1-z-w), [[1,1],[0,1]], 6),
        ('weighted', 1/(1-z-w), [[2,1],[1,1]], 3),
        ('signed_numerator', (1+2*z-w)/(1-z-w-z*w), [[1,1],[0,1]], 5),
        ('zero_row_and_column', (1+z+w)/(1-z-w), [[1,0,1],[0,0,0]], 3),
        ('all_zero', (2+z)/(1-z-w), [[0,0],[0,0]], 4),
        ('zero_at_origin', z*w/(1-z-w), [[1,1],[1,0]], 4),
    ]
    report = []
    for name, expr, matrix, bound in cases:
        out, groups = construct(expr, (z,w), matrix)
        variables = tuple(v for group in groups for v in group)
        original = rational_coefficients(expr, (z,w))
        transformed = rational_coefficients(out, variables)
        checks = 0
        for index in itertools.product(range(bound+1), repeat=len(groups)):
            old_index = tuple(sum(a*n for a,n in zip(row,index)) for row in matrix)
            new_index = tuple(index[j] for j,g in enumerate(groups) for _ in g)
            assert original(old_index) == transformed(new_index), (name,index)
            checks += 1
        report.append({'case':name, 'matrix':matrix, 'block_sizes':list(map(len,groups)),
                       'exact_coefficient_checks':checks, 'passed':True,
                       'rational_function':str(out)})
    result = {'all_passed':True, 'total_checks':sum(r['exact_coefficient_checks'] for r in report),
              'cases':report, 'proof_status':'finite checks supplement the symbolic proof'}
    Path(__file__).with_name('grouped_diagonal_checks.json').write_text(json.dumps(result,indent=2)+'\n')
    print(json.dumps({k:v for k,v in result.items() if k != 'cases'}))


if __name__ == '__main__':
    main()
\end{filecontents*}
\begin{filecontents*}[overwrite]{verify_grouped_diagonal.py}
"""Exact finite sanity checks of the coefficient-splitting construction.

These checks verify an implementation; they do not replace the symbolic proof.
"""

from functools import lru_cache
from itertools import product
import sympy as sp


def split(expression, old, new):
    left, right = new
    return sp.cancel((left * expression.subs(old, left)
                      - right * expression.subs(old, right)) / (left - right))


def construct(expression, variables, matrix):
    blocks = [[] for _ in matrix[0]]
    result = expression
    for row, old in enumerate(variables):
        row_variables = []
        for column, multiplicity in enumerate(matrix[row]):
            for copy in range(multiplicity):
                variable = sp.Symbol(f"x_{row}_{column}_{copy}")
                row_variables.append(variable)
                blocks[column].append(variable)
        if not row_variables:
            result = sp.cancel(result.subs(old, 0))
            continue
        current = old
        for index, variable in enumerate(row_variables[:-1]):
            remaining = sp.Symbol(f"temporary_{row}_{index}")
            result = split(result, current, (variable, remaining))
            current = remaining
        result = result.subs(current, row_variables[-1])
    for column, block in enumerate(blocks):
        if not block:
            dummy = sp.Symbol(f"dummy_{column}")
            block.append(dummy)
            result /= 1 - dummy
    return sp.cancel(result), blocks


def coefficient_function(expression, variables):
    numerator, denominator = map(sp.expand, sp.fraction(sp.cancel(expression)))
    p = dict(sp.Poly(numerator, *variables).terms())
    q = dict(sp.Poly(denominator, *variables).terms())
    zero = (0,) * len(variables)
    constant = q.pop(zero)
    assert constant != 0

    @lru_cache(None)
    def coefficient(index):
        if any(i < 0 for i in index):
            return sp.S.Zero
        value = p.get(index, sp.S.Zero)
        for exponent, scalar in q.items():
            value -= scalar * coefficient(tuple(i - e for i, e in zip(index, exponent)))
        return sp.cancel(value / constant)

    return coefficient


def main():
    z, w = sp.symbols("z w")
    rational_functions = [
        1 / (1 - z - w),
        (1 + z - 2*z*w) / (1 - 2*z + w - z*w - z*z),
    ]
    matrices = [
        ((1, 1), (1, 0)),
        ((2, 0), (0, 1)),
        ((1, 1), (0, 1)),
        ((0, 0), (1, 0)),
        ((0, 0), (0, 0)),
    ]
    checks = 0
    for original in rational_functions:
        source_coefficient = coefficient_function(original, (z, w))
        for matrix in matrices:
            transformed, blocks = construct(original, (z, w), matrix)
            variable_order = tuple(variable for block in blocks for variable in block)
            target_coefficient = coefficient_function(transformed, variable_order)
            for n in product(range(5), repeat=2):
                source_index = tuple(sum(a*b for a, b in zip(row, n)) for row in matrix)
                target_index = tuple(n[j] for j, block in enumerate(blocks) for _ in block)
                assert source_coefficient(source_index) == target_coefficient(target_index), (
                    original, matrix, n, transformed)
                checks += 1
    print(f"PASS: {checks} exact coefficient comparisons across "
          f"{len(rational_functions) * len(matrices)} constructions, "
          "including overlapping indices, weights, zero rows, and zero columns.")


if __name__ == "__main__":
    main()
\end{filecontents*}
\begin{filecontents*}[overwrite]{validate_problem4.py}
"""Exact arithmetic validation of the candidate CK4 bound; no numerics."""
from pathlib import Path
import json, time
import sympy as S
from sympy.polys.matrices import DomainMatrix

x,y=S.symbols('x y')

def polys(rx,ry):
    ux,dx=S.fraction(S.cancel(rx)); uy,dy=S.fraction(S.cancel(ry))
    D=S.lcm(dx,dy)
    A=S.cancel(rx*D); B=S.cancel(ry*D)
    den=S.ilcm(*[S.denom(c) for z in (A,B,D) for c in S.Poly(z,x,y).coeffs()])
    return tuple(S.Poly(S.expand(den*z),x,y,domain=S.ZZ) for z in (A,B,D))

def norm(p): return sum(abs(int(c)) for c in p.coeffs())
def deg(p,var): return max(0,int(p.degree(var))) if p else 0

def validate(name,rx,ry,s_override=None):
    start=time.monotonic()
    A,B,D=polys(rx,ry)
    t=max(deg(z,x) for z in (A,B,D)); s=max(deg(z,y) for z in (A,B,D))
    if s_override is not None: assert s_override>=s; s=s_override
    r=s+1; d=t*(2*s*s+2*s+1); a=d+(r-1)*t; b=(r-1)*s
    E=(d+r*t+1)*(r*s+1); np=(r+1)*(d+1); nu=(a+1)*(b+1)
    N=[S.Poly(1,x,y,domain=S.ZZ)]
    for i in range(r): N.append(D*N[-1].diff(x)+(A-i*D.diff(x))*N[-1])
    Ms=[N[i]*D**(r-i) for i in range(r+1)]
    columns=[]
    for i in range(r+1):
        for j in range(d+1): columns.append(Ms[i].mul(S.Poly(x**j,x,y)))
    C=B-(r-1)*D.diff(y)
    for j in range(a+1):
        for k in range(b+1):
            mon=S.Poly(x**j*y**k,x,y)
            columns.append(-(D*mon.diff(y)+C*mon))
    assert len(columns)==np+nu
    rows={}
    for j,col in enumerate(columns):
        for (ix,iy),c in col.terms():
            assert ix<=d+r*t and iy<=r*s
            if c: rows.setdefault(ix*(r*s+1)+iy,{})[j]=int(c)
    dm=DomainMatrix(rows,(E,np+nu),S.ZZ)
    ns=dm.nullspace().to_Matrix()
    candidates=[]
    for q in range(ns.rows):
        vv=[int(z) for z in ns.row(q)]
        if any(vv[:np]):
            g=S.igcd(*vv); vv=[z//g for z in vv]
            candidates.append(vv)
    assert candidates
    v=min(candidates,key=lambda q:(max(abs(z) for z in q),sum(bool(z) for z in q)))
    ps=[sum(v[i*(d+1)+j]*x**j for j in range(d+1)) for i in range(r+1)]
    U=sum(v[np+j*(b+1)+k]*x**j*y**k for j in range(a+1) for k in range(b+1))
    lhs=sum(S.Poly(ps[i],x,y)*Ms[i] for i in range(r+1))
    Up=S.Poly(U,x,y); rhs=D*Up.diff(y)+C*Up
    assert lhs==rhs
    alpha,beta,delta=map(norm,(A,B,D))
    product=1; hps=[]
    for i in range(r+1):
        hps.append(delta**(r-i)*product)
        product*=alpha+2*i*t*delta
    HP=max(hps); HU=max(1,beta+2*(r-1)*s*delta)
    assert all(norm(q)<=HP for q in columns[:np])
    assert all(norm(q)<=HU for q in columns[np:])
    K=min(max(HP,HU)**E,HP**np*HU**nu)
    assert max(abs(z) for z in v)<=K
    # Independently count the bounded certificate-only nullspace.
    ur={i:{j-np:c for j,c in row.items() if j>=np} for i,row in rows.items()}
    ur={i:row for i,row in ur.items() if row}
    udm=DomainMatrix(ur,(E,nu),S.ZZ)
    ukernel=nu-udm.rank()
    assert ukernel<=a+1
    result=dict(name=name,A=str(A.as_expr()),B=str(B.as_expr()),D=str(D.as_expr()),
      t=t,s=s,r=r,d=d,a=a,b=b,rows=E,columns=np+nu,rank=np+nu-ns.rows,
      certificate_kernel_dimension=ukernel,expected_dimension_surplus=np+nu-E-(a+1),
      actual_order=max(i for i,p in enumerate(ps) if p!=0),
      actual_degree=max(S.degree(p,x) for p in ps if p!=0),
      P_coefficients=[str(p) for p in ps],U=str(S.factor(U)),
      max_integer_coefficient=max(abs(z) for z in v),bound_bits=K.bit_length(),
      HP=HP,HU=HU,seconds=round(time.monotonic()-start,3),verified=True)
    print(json.dumps(result,default=int),flush=True)
    return result

CASES=[
 ('constant f=1',S.Integer(0),S.Integer(0),None),
 ('y-independent f=exp(x^2)',2*x,S.Integer(0),None),
 ('rational inverse f=1/(x+y), slack y bound',-1/(x+y),-1/(x+y),2),
 ('polynomial f=x+y',1/(x+y),1/(x+y),None),
 ('true hyperexponential f=exp(x*y)/(y^2+1)',y,x-2*y/(y*y+1),None),
 ('moving poles f=exp(x/(x+y))',y/(x+y)**2,-x/(x+y)**2,None),
 ('rational exponent f=(x+y)^(2/3)',S.Rational(2,3)/(x+y),S.Rational(2,3)/(x+y),None),
]

if __name__=='__main__':
    results=[validate(*case) for case in CASES]
    Path(__file__).with_name('problem4_validation.json').write_text(json.dumps(results,indent=2,default=int)+'\n')
\end{filecontents*}

\begin{document}
\maketitle
\begin{abstract}
We give two constructive statements associated with the problem list of
Chen and Kauers. First, multivariate binomial sums over a field of
characteristic zero are exactly grouped diagonals of rational power series.
The proof converts the monomial-section representation of Bostan, Lairez,
and Salvy into disjoint blocks of repeated indices by a classical
coefficient-splitting identity. A representation in $d$ variables with $s$
free indices can be converted using at most $ds$ variables when $d\ge1$.
Second, for a hyperexponential function with rational logarithmic
derivatives, we give explicit integer coefficient bounds for a polynomial
telescoper and its rational certificate. With logarithmic-derivative
bidegrees bounded by $(t,s)$, a possibly nonminimal telescoper can be chosen
with coefficient bit length $O(ts^3(h+\log(ts+2)))$, where $h$ measures the
input coefficient norms and $t,s\ge1$. All exceptional cases in the
construction are treated. The proofs use established algebraic techniques;
the manuscript does not certify historical priority or a new foundational
method. Exact supplementary checks are embedded in the single TeX source.
\end{abstract}

\section{Questions, results, and attribution}
This note uses the original arXiv numbering in Chen--Kauers
\cite{CK2016}. Problem~12 asks for a characterization of multivariate
binomial sums by coefficients of rational power series with each free
index repeated in a separate block. Problem~4 asks for bounds on the bit
size of hyperexponential telescopers. Their journal numbers are,
respectively, 6.3 and 3.2.

Theorem~\ref{thm:gd-characterization} proves the former equivalence for
sequences on $\mathbb N^s$, the index domain implicit in ordinary power
series. Theorem~\ref{he:main} gives explicit bit bounds for every
hyperexponential input whose logarithmic derivatives belong to
$\mathbb Q(x,y)$, under a stated finite arithmetic representation.
The latter theorem asserts existence at bounded order and degree, without
requiring the order to be minimal.

\paragraph{Dependence on earlier work.}
The substantial rational monomial-section construction in the first
argument is due to Bostan--Lairez--Salvy \cite{BLS2017}. Its remaining
conversion is proved here using the classical Hankel-array identity
recorded, for example, by Barry \cite[p.~5]{Barry2016}.
The second argument uses polynomial ansatz, dimension counting, and
Hadamard bounds. These techniques have established precedents in
creative telescoping \cite{CK2012,KY2015}.
The resulting statements should be assessed as explicit constructive
derivations, not evidence that their underlying tools originated here.

\paragraph{Priority and validation.}
A search of the primary sources considered in preparing this note did
not locate an explicit identical full result. That search cannot establish
the absence of prior work, and no claim of first resolution is made.
The two proofs were checked independently during preparation; the
additional computations are exact rational or integer checks, not formal
proof-assistant verification or external peer review.

The distinction between algebraic complexity and coefficient bit size
matters for the second question. Van der Hoeven's general reduction
framework \cite[Section~8.1 and Theorem~9.8]{vdH2018} measures field
operations and coefficient counts. Its stated polynomial complexity does
not by itself provide the explicit integer coefficient bounds proved below.
This observation does not imply that no other source contains such bounds.

\section{Grouped diagonals and multivariate binomial sums}
\label{sec:gd}

Chen and Kauers ask whether the multivariate binomial sums of
Bostan, Lairez, and Salvy admit rational representations in which each free
index is repeated in its own block of variables
\cite[Problem~12]{CK2016}.  The journal version labels this question
Problem~6.3.  We give an affirmative proof by converting a nonnegative
monomial section into a grouped diagonal.  The input representation is due
to Bostan, Lairez, and Salvy; the conversion uses an elementary
divided-difference operation on coefficient indices.

Throughout this section, $K$ is a field of characteristic zero and
$\mathbb N=\{0,1,2,\ldots\}$.  A \emph{rational power series} is an element
of $K(\boldsymbol z)\cap K[[\boldsymbol z]]$ represented as $P/Q$ with
$P,Q\in K[\boldsymbol z]$ and $Q(\boldsymbol 0)\ne0$.
We use the binomial-sum class of \cite[Definition~1.1]{BLS2017}: it is
generated by Kronecker deltas, nonzero geometric sequences, and binomial
coefficients, under addition, multiplication, affine integer changes of
indices, and directed indefinite summation.  A sequence on $\mathbb N^s$
belongs to this class when its zero extension to $\mathbb Z^s$ does.
This convention is legitimate because the indicator
$H(n)=\boldsymbol 1_{n\ge0}$ is itself a binomial sum.

The nonnegative index set is the one implicit in the rational-power-series
formulation of the question.  A two-sided sequence on $\mathbb Z^s$ has
additional negative-index data, which a single ordinary power series does
not encode.  We do not reinterpret those data as coefficients with
negative exponents.

\begin{theorem}[Grouped-diagonal characterization]
\label{thm:gd-characterization}
Let $s\ge1$ and let $b\colon\mathbb N^s\longrightarrow K$.
The following conditions are equivalent.
\begin{enumerate}
\item The sequence $b$ is a binomial sum.
\item There exist positive integers $r_1,\ldots,r_s$ and a rational
power series $U$ in variables $x_{j,\ell}$,
$1\le j\le s$, $1\le\ell\le r_j$, such that, for every
$\boldsymbol n\in\mathbb N^s$,
\begin{equation}
 b(\boldsymbol n)=
 \left[\prod_{j=1}^s\prod_{\ell=1}^{r_j}x_{j,\ell}^{n_j}\right]U.
 \label{eq:gd-target}
\end{equation}
\end{enumerate}
Both implications are constructive.
\end{theorem}

\subsection{The existing monomial-section representation}

We first state precisely the result from which the conversion starts.
It is the intermediate construction in the proof of
\cite[Proposition~3.10 and Section~3.1]{BLS2017}.

\begin{proposition}[Bostan--Lairez--Salvy monomial-section input]
\label{prop:gd-bls}
For every binomial sum $b\colon\mathbb N^s\to K$, one can construct a
rational power series
\[
 T(\boldsymbol z)=\sum_{\boldsymbol m\in\mathbb N^d}
 c_{\boldsymbol m}\boldsymbol z^{\boldsymbol m}
\]
and a matrix $A\in\mathbb N^{d\times s}$ such that
\begin{equation}
 b(\boldsymbol n)=c_{A\boldsymbol n}
 \qquad(\boldsymbol n\in\mathbb N^s).
 \label{eq:gd-monomial-section}
\end{equation}
\end{proposition}

For clarity, the nonnegativity and origin regularity in this statement
are substantial parts of the existing construction.  The initial
representation is a constant term of a rational function in an ordered
Laurent-series field.  The nonnegative integral shear of
\cite[Lemma~3.11]{BLS2017} turns its denominator into
$C\boldsymbol z^{\boldsymbol\mu}(1+G)$, with
$C\ne0$, $\boldsymbol\mu\in\mathbb N^d$, and $G(\boldsymbol0)=0$.
The shear is injective and respects the chosen Laurent-series expansion;
the target monomials therefore become $\boldsymbol z^{A\boldsymbol n}$
with $A\ge0$.  At this stage the representation is
\[
 b(\boldsymbol n)=
 [\boldsymbol z^{A\boldsymbol n}]
 \boldsymbol z^{-\boldsymbol\mu}S(\boldsymbol z),
 \qquad S\text{ rational and regular at }\boldsymbol0.
\]
If $\mu_i>0$, replace $S$ by
$(S-S|_{z_i=0})/z_i$ and reduce $\mu_i$ by one.  This preserves every
target coefficient: the discarded term has exponent $-\mu_i<0$ in
$z_i$, whereas $(A\boldsymbol n)_i\ge0$.
Each replacement is rational and regular, and exactly
$\sum_i\mu_i$ such replacements suffice.  This gives
\eqref{eq:gd-monomial-section}.  This explanation also shows why the
one-index argument in Section~3.1 applies to the several target
monomials in Proposition~3.10.

The columns of $A$ can overlap in their nonzero rows.  Consequently,
\eqref{eq:gd-monomial-section} is not already the repeated-block
representation \eqref{eq:gd-target}.  The following operation removes
that overlap.

\subsection{Splitting coefficient indices}

The identity below is the familiar generating function of a Hankel
array; it appears explicitly, for example, in \cite[p.~5]{Barry2016}.
We include its proof with auxiliary variables and origin regularity,
which are needed for the repeated construction.

\begin{lemma}[Binary splitting]
\label{lem:gd-split}
Let $F(z,\boldsymbol u)$ be rational and regular at the origin, with
\[
 F(z,\boldsymbol u)=\sum_{n\ge0}f_n(\boldsymbol u)z^n.
\]
Then
\begin{equation}
 \mathcal S_zF(x,y,\boldsymbol u)
 =\frac{xF(x,\boldsymbol u)-yF(y,\boldsymbol u)}{x-y}
 =\sum_{p,q\ge0}f_{p+q}(\boldsymbol u)x^py^q
 \label{eq:gd-split}
\end{equation}
is also rational and regular at the origin.
\end{lemma}

\begin{proof}
The coefficient identity follows from
\[
 \frac{x^{n+1}-y^{n+1}}{x-y}
 =\sum_{p+q=n}x^py^q.
\]
To check regularity rather than just rationality, write $F=P/Q$ with
$Q(0,\boldsymbol0)\ne0$.  Over the denominator
$Q(x,\boldsymbol u)Q(y,\boldsymbol u)$ the numerator before division
by $x-y$ is
\[
 xP(x,\boldsymbol u)Q(y,\boldsymbol u)
 -yP(y,\boldsymbol u)Q(x,\boldsymbol u).
\]
This polynomial vanishes at $x=y$, hence is divisible by $x-y$.
After cancellation the denominator is a unit at the origin.  Thus the
apparent pole along $x=y$ is removable by polynomial division.
The identity is coefficientwise and requires no interchange of
infinite sums with infinitely many contributors to a coefficient.
\end{proof}

Repeated splitting replaces a coefficient index $m$ by a sum of any
prescribed positive number $q$ of independent indices.  With $q=1$
nothing is changed, and with $q=0$ the appropriate operation is
specialization of the original variable to zero.

\begin{lemma}[Monomial sections as grouped diagonals]
\label{lem:gd-section}
Let $T=\sum_{\boldsymbol m\in\mathbb N^d}
c_{\boldsymbol m}\boldsymbol z^{\boldsymbol m}$ be a rational power
series, and let $A\in\mathbb N^{d\times s}$.  The sequence
$c_{A\boldsymbol n}$ has a representation of the form
\eqref{eq:gd-target}.  Such a representation can be constructed using
\[
 \sum_{i=1}^d\sum_{j=1}^s A_{ij}
 +\#\{j:A_{ij}=0\text{ for every }i\}
\]
variables.
\end{lemma}

\begin{proof}
For each row $i$, put $q_i=\sum_j A_{ij}$.  If $q_i=0$, specialize
$z_i=0$.  Otherwise use Lemma~\ref{lem:gd-split} to replace $z_i$ by
$q_i$ independent variables, labelled
\[
 x_{i,j,\ell},\qquad 1\le j\le s,\quad1\le\ell\le A_{ij}.
\]
Carry this out for all rows, obtaining a regular rational series $V$.
For arbitrary nonnegative exponents $k_{i,j,\ell}$ its coefficients are
\begin{equation}
 \left[\prod_{i,j,\ell}x_{i,j,\ell}^{k_{i,j,\ell}}\right]V
 =c_{\boldsymbol m},\qquad
 m_i=\sum_{j,\ell}k_{i,j,\ell}.
 \label{eq:gd-split-array}
\end{equation}
In particular, requiring $k_{i,j,\ell}=n_j$ gives
$m_i=\sum_j A_{ij}n_j$.  All variables with middle label $j$ form
the $j$th block.  If column $j$ is zero, multiply $V$ by
$(1-v_j)^{-1}$ in a fresh variable and use $\{v_j\}$ as that block.
Its coefficient is one at every nonnegative index, so the desired
sequence is unchanged.  All blocks are now nonempty.  Reordering the
variables gives \eqref{eq:gd-target}, with the asserted variable count.
The constant case $d=0$ is handled by
$U=T\prod_{j=1}^s(1-v_j)^{-1}$.
\end{proof}

\subsection{Proof of the characterization}

\begin{proof}[Proof of Theorem~\ref{thm:gd-characterization}]
For the forward implication, apply Proposition~\ref{prop:gd-bls} and
then Lemma~\ref{lem:gd-section}.

For the converse we give an explicit finite-sum argument.  By scaling
the denominator and separating the monomials of the numerator, it
suffices to consider
\begin{equation}
 R(\boldsymbol z)=
 \frac{\boldsymbol z^{\boldsymbol\beta}}
 {1-\sum_{\ell=1}^e a_\ell
                    \boldsymbol z^{\boldsymbol\alpha_\ell}},
 \label{eq:gd-converse-rational}
\end{equation}
where $\boldsymbol\beta\in\mathbb N^r$,
$\boldsymbol\alpha_\ell\in\mathbb N^r\setminus\{\boldsymbol0\}$,
and $a_\ell\in K\setminus\{0\}$.
Its coefficient at $\boldsymbol z^{\boldsymbol m}$ equals
\begin{equation}
 \sum_{\substack{\boldsymbol k\in\mathbb N^e\\
 \boldsymbol\beta+\sum_\ell k_\ell\boldsymbol\alpha_\ell
                         =\boldsymbol m}}
 \binom{k_1+\cdots+k_e}{k_1,\ldots,k_e}
 \prod_{\ell=1}^e a_\ell^{k_\ell}.
 \label{eq:gd-converse-sum}
\end{equation}
Every contributing tuple satisfies
$\sum_\ell k_\ell\le |\boldsymbol m|_1$, because each
$|\boldsymbol\alpha_\ell|_1\ge1$.
Thus \eqref{eq:gd-converse-sum} can be written as nested finite sums
with bounds $0\le k_\ell\le|\boldsymbol m|_1$, inserting one
Kronecker delta for each coordinate of the displayed equality.
These bounds are affine in $\boldsymbol m$.  The multinomial factor is
a product of binomial coefficients, and the remaining factors are
permitted geometric sequences and deltas.  The defining closure rules
therefore make the entire coefficient array a binomial sum on
$\mathbb N^r$.  Multiplication by $\prod_i H(m_i)$ supplies its zero
extension.  When $e=0$, the coefficient array is simply a product of
deltas, so the same conclusion holds.  Finally, the repeated-block
map from $\boldsymbol n$ to $\boldsymbol m$ is an affine integer
substitution, under which binomial sums are closed.
This proves the converse and also its effectivity.
\end{proof}

\subsection{Example and a smaller variable bound}

Take $T(z,w)=(1-z-w)^{-1}$ and
\[
 A=\begin{pmatrix}1&1\\1&0\end{pmatrix}.
\]
The corresponding monomial section is
$b(m,n)=\binom{2m+n}{m}$.  Splitting the $z$ index gives
\begin{equation}
 U(a,b,c)=\frac{1-c}{(1-a-c)(1-b-c)},\qquad
 [a^mc^mb^n]U=\binom{2m+n}{m}.
 \label{eq:gd-example}
\end{equation}
Indeed, $[a^mb^n]U=(1-c)^{-m-n-1}$.  Thus the shared index $m+n$
in the original $z$ coordinate becomes two independent variables,
subsequently assigned to different blocks.

The copies-only construction above is sufficient for the theorem.
For large entries of $A$, ordinary rational multisections reduce its
number of variables.

\begin{proposition}[A dimension bound]
\label{prop:gd-dimension}
If $d\ge1$, the representation in Lemma~\ref{lem:gd-section} can be
constructed with at most $ds$ variables, independently of the sizes
of the entries of $A$.
\end{proposition}

\begin{proof}
Split row $i$ into one variable $x_{ij}$ for each nonzero entry
$A_{ij}$, rather than into $\sum_j A_{ij}$ variables.
The resulting coefficient array is
$c_{(\sum_j k_{1j},\ldots,\sum_j k_{dj})}$.
In coordinate $x_{ij}$ take the $A_{ij}$-multisection: retain the
coefficients whose exponent is divisible by $A_{ij}$ and divide that
exponent by $A_{ij}$.

To justify this operation over $K$, let $F(x,\boldsymbol u)=P/Q$ be
regular at the origin and let $a\ge1$.  In a finite splitting field
for the $a$th roots of unity form
\[
 W(x,\boldsymbol u)=\frac1a\sum_{\zeta^a=1}F(\zeta x,\boldsymbol u).
\]
Its series contains exactly the terms with $x$ exponent divisible by
$a$.  A denominator for $W$ is
$D=\prod_{\zeta^a=1}Q(\zeta x,\boldsymbol u)$, whose constant term
is nonzero.  Both $D$ and the polynomial numerator $DW$ are invariant
under $x\mapsto\zeta x$ and under the Galois group of the splitting
field.  Hence both belong to $K[x^a,\boldsymbol u]$.  Replacing
$x^a$ by a new variable gives a rational power series over $K$ with
coefficient $[x^{an}]F$ at exponent $n$.

Apply these multisections independently.  The coefficient array becomes
\[
 c_{(\sum_j A_{1j}k_{1j},\ldots,\sum_j A_{dj}k_{dj})}.
\]
Its grouped diagonal is the required sequence.
Add one geometric dummy variable for each zero column, as before.
If there are $q$ zero columns, the variable count is at most
$d(s-q)+q\le ds$.
\end{proof}

\begin{remark}[Scope of the contribution]
\label{rem:gd-priority}
The proof uses the established Bostan--Lairez--Salvy representation
and classical divided differences and multisections.  Its additional
step is the explicit conversion from overlapping nonnegative monomial
sections to disjoint repeated-index blocks.  This supplies a proof of
the stated characterization; it does not, by itself, establish
publication priority or a historically new foundational principle.
\end{remark}

\section{An explicit arithmetic bound for hyperexponential telescopers}
\label{he:section}

The original arXiv Problem~4 of Chen and Kauers~\cite{CK2016}, numbered
Problem~3.2 in the journal version, asks for bounds on the bit size of
telescopers for hyperexponential terms. We give an explicit bound in the
rational logarithmic-derivative input model. The proof is self-contained,
but its methodology is classical: polynomial ansatz and dimension counting,
as in the order--degree analysis of Chen and Kauers~\cite{CK2012}, followed by
integer-kernel bounds using determinants. Arithmetic size analysis for
hypergeometric telescopers was already undertaken by Kauers and Yen~\cite{KY2015}.
No claim of a historically new method or established priority is made here.

Let $\mathcal F$ be a differential field extension of $\mathbb Q(x,y)$
with commuting derivations $\partial_x,\partial_y$, and let
$0\ne f\in\mathcal F$. Suppose we are given
$A,B,D\in\mathbb Z[x,y]$, $D\ne0$, such that
\begin{equation}\label{he:input}
 \frac{\partial_x f}{f}=\frac AD,
 \qquad \frac{\partial_y f}{f}=\frac BD.
\end{equation}
Every pair of logarithmic derivatives in $\mathbb Q(x,y)$ admits this
representation by clearing polynomial and integer denominators.
For $F\in\mathbb Z[x,y]$, write $\|F\|_1$ for the sum of the absolute
values of its coefficients. The coefficient magnitude bit length of a
nonzero integer polynomial operator is
$1+\lfloor\log_2\max |c|\rfloor$, where $c$ runs over its integer
coefficients; sign bits are counted separately.

\begin{theorem}\label{he:main}
Let $t,s\ge0$ be integer upper bounds on the respective $x$- and
$y$-degrees of $A,B,D$. Let $r,d$ be nonnegative integers satisfying
\begin{equation}\label{he:admissible}
 r>s,\qquad
 (r-s)(d+1)>t\bigl((2r-1)s+r\bigr).
\end{equation}
Put
\begin{align*}
 a&=d+(r-1)t,& b&=(r-1)s,\\
 n_P&=(r+1)(d+1),&n_U&=(a+1)(b+1),\\
 E&=(d+rt+1)(rs+1).
\end{align*}
Define $\alpha=\|A\|_1$, $\beta=\|B\|_1$, $\delta=\|D\|_1\ge1$, and
\begin{align}
 H_P&=\max_{0\le i\le r}
  \left\{\delta^{r-i}\prod_{j=0}^{i-1}(\alpha+2jt\delta)\right\},
       \label{he:hp}\\
 H_U&=\max\{1,\beta+2(r-1)s\delta\},\qquad
 H=\max\{H_P,H_U\},\label{he:hu}\\
 K&=\min\{H^E,H_P^{n_P}H_U^{n_U}\}.\label{he:K}
\end{align}
An empty product equals one. There exist $p_0,\ldots,p_r\in\mathbb Z[x]$,
not all zero, and $U\in\mathbb Z[x,y]$, with
\[
 \deg_xp_i\le d,\qquad \deg_xU\le a,\qquad \deg_yU\le b,
\]
such that
\begin{equation}\label{he:identity}
 \sum_{i=0}^r p_i(x)\partial_x^i f
   =\partial_y\left(\frac{U}{D^{r-1}}f\right).
\end{equation}
Every integer coefficient of the $p_i$ and $U$ has absolute value at most
$K$. In particular, the nonzero telescoper
$P=\sum_{i=0}^rp_i(x)\partial_x^i$ has order at most $r$, degree at most
$d$, and coefficient magnitude bit length at most
\begin{equation}\label{he:bits}
 1+\left\lceil\min\left\{
 E\log_2H,
 n_P\log_2H_P+n_U\log_2H_U\right\}\right\rceil.
\end{equation}
Its dense representation has at most $n_P$ coefficients.
\end{theorem}

\begin{proof}
Define $N_0=1$ and, for $i\ge0$,
\begin{equation}\label{he:recurrence}
 N_{i+1}=D\partial_xN_i+(A-i\partial_xD)N_i.
\end{equation}
Induction using~\eqref{he:input} gives
$\partial_x^if/f=N_i/D^i$ and
$\deg_xN_i\le it$, $\deg_yN_i\le is$. Thus
$M_i=N_iD^{r-i}$ has bidegree at most $(rt,rs)$.

Make a polynomial ansatz for the $p_i$ and $U$ with the bounds in the
statement. After multiplication by $D^r/f$, identity~\eqref{he:identity}
becomes
\begin{equation}\label{he:polysystem}
 \sum_{i=0}^r p_iM_i
 -D\partial_yU-\bigl(B-(r-1)\partial_yD\bigr)U=0.
\end{equation}
Its bidegree is at most $(d+rt,rs)$. Equating coefficients gives a
homogeneous integer matrix $M$ with $E$ rows and $n_P+n_U$ columns.
Zero rows can be included.

We must ensure that a kernel vector has a nonzero $p$-component.
The certificate-only equation is
\begin{equation}\label{he:purecert}
 D\partial_yU+\bigl(B-(r-1)\partial_yD\bigr)U=0.
\end{equation}
Over $\mathbb Q(x)$ its polynomial solution space has dimension at most
one: the ratio of two nonzero solutions has zero $y$-derivative and
therefore belongs to $\mathbb Q(x)$. If a bounded nonzero solution $U_0$
exists, choose one of its nonzero $y$-coefficients. Extracting that same
coefficient from any bounded solution is injective, because every solution
is $q(x)U_0$. Its image lies in $\mathbb Q[x]_{\le a}$. Hence the
$\mathbb Q$-dimension of the bounded certificate-only kernel is at most
$a+1$. Meanwhile
\begin{align*}
 \dim_{\mathbb Q}\ker M-(a+1)
 &\ge n_P+n_U-E-(a+1)\\
 &=(r-s)(d+1)-t\bigl((2r-1)s+r\bigr)>0.
\end{align*}
Thus some kernel vector has a nonzero $p$-component. This argument includes
zero logarithmic derivatives and the cases $s=0$ or $t=0$.

Next we bound column norms. Since
$\|\partial_xN_i\|_1\le it\|N_i\|_1$ and
$\|\partial_xD\|_1\le t\delta$, recurrence~\eqref{he:recurrence} yields
\[
 \|N_{i+1}\|_1\le(\alpha+2it\delta)\|N_i\|_1.
\]
Consequently $\|M_i\|_1\le H_P$. Each column belonging to a coefficient
of $p_i$ is a monomial translate of $M_i$, so its coefficient-vector
$\ell_1$ norm is at most $H_P$. For a coefficient of $x^jy^k$ in $U$,
where $0\le k\le b$, the corresponding column is, up to sign and an
$x$-monomial translate, the coefficient vector of
\[
 kDy^{k-1}+\bigl(B-(r-1)\partial_yD\bigr)y^k.
\]
Its norm is at most
$k\delta+\beta+(r-1)s\delta\le H_U$.

Let $\rho=\operatorname{rank}M$. If $\rho=0$, a unit vector in any
$p$-coordinate proves the assertion. Otherwise choose $\rho$ independent
columns and $\rho$ rows whose intersection $C$ is nonsingular. For each
nonpivot column $k$, give its free coordinate the value $\det C$, set
the other free coordinates to zero, and set the pivot coordinates equal
to $-\operatorname{adj}(C)M[I,k]$. These integer vectors solve all rows:
the pivot columns span the full column space, and restriction to the
selected rows is injective on that space. They form a basis of the
rational kernel.

Every coordinate of such a vector is zero or a $\rho$-by-$\rho$ minor
of $M$. Hadamard's inequality, using that a column's $\ell_2$ norm is at
most its $\ell_1$ norm, bounds each minor by $H^\rho\le H^E$.
Since a minor uses distinct columns, it is also bounded by
$H_P^{n_P}H_U^{n_U}$; both bases are at least one. At least one basis
vector has a nonzero $p$-component, because the full kernel does. That
vector proves~\eqref{he:identity} and~\eqref{he:K};
the bit estimate follows.
\end{proof}

\begin{proposition}\label{he:explicit}
Admissible choices in Theorem~\ref{he:main} are
\[
 (r,d)=\bigl(s+1,t(2s^2+2s+1)\bigr)
 \quad\hbox{and}\quad
 (r,d)=\bigl(2s+1,t(4s+1)\bigr).
\]
More generally, for any integer $r>s$ one may take
\[
 d=\left\lfloor\frac{t((2r-1)s+r)}{r-s}\right\rfloor.
\]
If $t,s\ge1$ and
$h=\log_2\max\{1,\|A\|_1,\|B\|_1,\|D\|_1\}$, the second uniform
choice yields a telescoper with coefficient bit length
\[
 O\bigl(ts^3(h+\log(ts+2))\bigr).
\]
\end{proposition}
\begin{proof}
For the first choice the difference between the two sides of the strict
inequality in~\eqref{he:admissible} equals one. For the second choice it
equals $2ts+s+1$. For the floor choice, $d+1$ is strictly greater than
the displayed quotient.
For the asymptotic bound, $r=O(s)$, $d=O(ts)$,
$n_P=O(ts^2)$, and $n_U=O(ts^3)$. Formula~\eqref{he:hp} gives
$\log_2H_P=O(s(h+\log(ts+2)))$, while
$\log_2H_U=O(h+\log(s+2))$. Substitute these bounds into the second
term of~\eqref{he:bits}.
\end{proof}

\begin{remark}\label{he:scope}
The assertion concerns existence of a bounded, possibly nonminimal-order
telescoper. It neither bounds every scalar multiple nor promises minimum
order. Common integer content can safely be divided out of $P$, with a
rational rescaling of the certificate; no height claim is made for
division by an arbitrary common polynomial factor. Rational exponents
are included when the logarithmic derivatives lie in $\mathbb Q(x,y)$.
An arithmetic bound for arbitrary irrational constants needs an explicit
encoding and is not part of this statement. If the input polynomial
coefficients have magnitude bit length at most $c$, then
$h\le c+\log_2((t+1)(s+1))$.
\end{remark}

\paragraph{Exact checks.}
The polynomial matrices in the proof were constructed over $\mathbb Z$
for seven examples. Exact nullspaces were computed with SymPy~1.14.0;
the resulting polynomial identity, the certificate-only kernel bound,
every column-norm estimate, and the final integer-height bound were
checked exactly. Table~\ref{he:validation} lists the actual operators
selected by the computation, not claimed minimal operators. Here
$\kappa$ denotes the dimension of the bounded certificate-only kernel.
All examples use the first uniform choice in
Proposition~\ref{he:explicit}; the third uses a deliberately nonsharp
$y$-degree bound $s=2$ to include a nonzero certificate-only kernel.

\begin{table}[htbp]
\centering
\small
\begin{tabular}{lrrrl}
\hline
$f$ & $r$ & $d$ & $\kappa$ & Computed $P$\\
\hline
$1$ &1&0&1&$\partial_x$\\
$e^{x^2}$ &1&1&2&$\partial_x-2x$\\
$(x+y)^{-1}$ &3&13&13&$\partial_x^3$\\
$x+y$ &2&5&7&$\partial_x^2$\\
$e^{xy}/(y^2+1)$ &4&25&0&$x(\partial_x^2+1)$\\
$e^{x/(x+y)}$ &3&26&0&$x\partial_x^3+\partial_x^2$\\
$(x+y)^{2/3}$ &2&5&0&$\partial_x^2$\\
\hline
\end{tabular}
\caption{Exact validation of the polynomial ansatz and arithmetic bounds.}
\label{he:validation}
\end{table}

For example, with $f=e^{xy}/(y^2+1)$ we used
$A=y^3+y$, $B=x(y^2+1)-2y$, $D=y^2+1$.
The selected solution has $U=(y^2+1)^4$, so its certificate is
$(y^2+1)f=e^{xy}$, and the identity is
$x(\partial_x^2+1)f=\partial_y(e^{xy})$.
For the moving-pole example the certificate numerator is
$U=y^2(2x+y)$ with $D=(x+y)^2$ and $r=3$.
These checks supplement the proof; they do not establish novelty.

\paragraph{Interpretation.}
The theorem supplies an explicit arithmetic bound for all inputs in
model~\eqref{he:input}, addressing the literal bound request of
Problem~4 in that model. The bound is deliberately coarse and no matching
lower bound is asserted. Its independent derivation and verification do
not prove absence of an identical or stronger result in the literature.
In particular, the stronger requirements of a historically new principle
and an established first resolution remain unverified.

\section{Reproducible checks and remaining limitations}
The mathematical arguments above apply to all permitted inputs and
indices. The following finite checks are supplementary regression checks
for the constructions and for common implementation errors.

Two independently written coefficient-checking programs verified
465 exact coefficient comparisons in 16 rational constructions.
The cases include overlapping and weighted index maps, zero rows,
zero columns, an all-zero map, signed numerator coefficients, and
vanishing constant terms. Each program constructs and cancels the
rational divided differences before independently extracting coefficients.
The denominator is checked to be nonzero at the origin.

The arithmetic-bound checker constructs seven exact integer matrices,
including a $390\times420$ example and a $231\times263$ example.
It checks an actual telescoper-certificate identity, the dimension of the
certificate-only kernel, all column-norm bounds, and the stated integer
coefficient bound. It also tests vanishing logarithmic derivatives,
rational input functions, a rational exponent, and a moving pole.

\paragraph{Single-source reproduction.}
This TeX file contains three complete Python programs in
\texttt{filecontents*} blocks. Compiling it with a current pdfLaTeX
installation writes the programs alongside the PDF. With Python~3 and
SymPy installed, execute
\begin{verbatim}
python verify_grouped_diagonals.py
python verify_grouped_diagonal.py
python validate_problem4.py
\end{verbatim}
The first and third programs write JSON reports. No floating-point
tolerance is used to accept a mathematical identity. The manuscript
contains its own bibliography and needs no external TeX inputs or figures.

\paragraph{What is not established.}
The results do not prove that either construction is the first in the
literature. They do not introduce a historically new basic principle.
The arithmetic theorem does not assert an optimal bound or minimal-order
coefficient bound, and it does not assign a bit size to unencoded complex
constants. No further problem from the eighteen-item list is asserted to
be resolved by this note.

\begin{thebibliography}{9}
\bibitem{CK2016}
S. Chen and M. Kauers,
\emph{Some open problems related to creative telescoping},
Journal of Systems Science and Complexity \textbf{30} (2017), 154--172.
Original numbered list: \href{https://arxiv.org/abs/1609.03768}{arXiv:1609.03768v1} (2016).

\bibitem{BLS2017}
A. Bostan, P. Lairez, and B. Salvy,
\emph{Multiple binomial sums},
Journal of Symbolic Computation \textbf{80} (2017), 351--386.
\href{https://arxiv.org/abs/1510.07487}{arXiv:1510.07487v2} (2016).

\bibitem{Barry2016}
P. Barry,
\emph{Jacobsthal Decompositions of Pascal's Triangle, Ternary Trees,
and Alternating Sign Matrices},
Journal of Integer Sequences \textbf{19} (2016), Article 16.3.5.
\url{https://cs.uwaterloo.ca/journals/JIS/VOL19/Barry/barry321.pdf}.

\bibitem{CK2012}
S. Chen and M. Kauers,
\emph{Trading order for degree in creative telescoping},
Journal of Symbolic Computation \textbf{47} (2012), 968--995.
\href{https://arxiv.org/abs/1108.4508}{arXiv:1108.4508}.

\bibitem{KY2015}
M. Kauers and L. Yen,
\emph{On the length of integers in telescopers for proper hypergeometric terms},
Journal of Symbolic Computation \textbf{66} (2015), 21--33.
\href{https://arxiv.org/abs/1311.3720}{arXiv:1311.3720}.

\bibitem{vdH2018}
J. van der Hoeven,
\emph{Creative telescoping using reductions}, draft, 8 June 2018.
\url{https://www.texmacs.org/joris/redtel/redtel.pdf}.
\end{thebibliography}
\end{document}
